\chapter{The Second Variation Is the Single Invariant}
\label{chap:second-variation-single-invariant}

This is the chapter the whole second part has been built to reach, and it earns being stated plainly.
The ninth chapter cleared a stage by sending the first variation to zero. The eleventh found, at the
interior node of a spline, a mixed coupling term and pointedly carried it whole, calling it the seed of
something. The twelfth built an energy whose one virtue was that it detects zero, and earned its
positive-definiteness at a boundary. Each of those chapters was circling the same object without
naming it. This chapter names it. There is a single invariant at the center of the construction --- the
book is named for it --- and it is the \emph{second variation}: the second-order response of an action
to the variation of a path. The theorem of this chapter is that this one object is reached by three
different roads that all arrive at the same place, and that its arrival is what the rest of the book
will read.

\section{The second variation, named}
\label{se:the-second-variation-named}

Begin with the mixed coupling. The eleventh chapter showed that varying the interior node of a cubic
path splits the change in the action into three pieces --- a left residual, a right residual, and a
mixed term in which the two halves' variations multiply together --- and it left that mixed term
untouched, refusing to approximate it away. The construction now gives it its name: the mixed coupling
\emph{is} the second variation. The definition is exactly that direct; there is no new object built,
only a name placed on the term that was carried whole for two chapters waiting for it. And the name is
earned by what the term does when the first variation is gone. Drive the left and right residuals to
zero --- make the path stationary in each half, as a solved variational problem does --- and the entire
coupled difference collapses to the mixed term alone. What survives once the first variation has
vanished is exactly the second variation, and nothing else remains to compete with it. The first
variation lived in the left and right; the second variation is what is left when they are spent.

It is worth watching the collapse happen, because it is the mechanism that makes the second variation
visible. Write out the coupled difference at a path that is not yet stationary: three terms, the left
residual large, the right residual large, the mixed term a modest product beneath them. Now solve the
variational problem --- find the stationary path, the one whose first variation vanishes in each half.
The left residual goes to zero; the right residual goes to zero; and the mixed term, which depended on
neither of them, is untouched. Where there were three terms there is now one, and it is the mixed
coupling standing alone. The second variation is what the construction funges out of the coupled
difference by the single characteristic of surviving the first variation's vanishing: of the three
contributions, it is the one that does not depend on the path being non-stationary, so it is the one
that remains when stationarity is imposed. That funging --- pooling away the two first-order residuals
and keeping the second-order remainder --- is the operation that isolates the invariant.

This is why the first variation had to be cleared before the second could be seen. As long as the path
was not yet stationary, the change in the action was dominated by the first-order residuals, and the
mixed term was a small correction buried beneath them. At a stationary path the first-order residuals
are zero by construction --- that is what stationary means --- and the mixed term, no longer buried,
becomes the leading description of how the action behaves. The second variation is the action's
behavior at a stationary path, the structure that governs a path once it has been made flat. To get
to it, the construction had to spend Part II flattening the first variation; the reward for that labor
is that one term, the mixed coupling, is left standing alone, ready to be measured.

\section{The unique quadratic remainder}
\label{se:the-unique-quadratic-remainder}

The second road approaches the same object from the side of remainders. When the first variation of an
action is taken and found to vanish, the action has not been fully described --- it has only been
described to first order, and there is a remainder, the part the linear approximation missed. The
construction proves that this remainder is quadratic and, more, that it is \emph{unique}: there is one
and only one second-order term that correctly accounts for what the first variation left behind. The
second variation is not a choice among possible second-order corrections; it is the determined one, the
remainder that is forced by the requirement of describing the action correctly past first order.

Uniqueness is what makes it an invariant rather than an artifact. A quantity that depended on how the
remainder was carved up, or on which second-order term one chose to write, would be a feature of the
description, not of the action. Because the second-order remainder is unique, it is a fact about the
action itself, the same no matter how the bookkeeping is arranged --- which is precisely the test the
construction has applied to every claimed invariant since the Galileo experiment of the first chapter:
a thing is physical exactly when it survives every admissible way of describing it. The second
variation survives. The mixed coupling of the first road and the unique quadratic remainder of this one
are, the construction proves, the same term, so the object has now been reached twice and found
identical both times.

It is worth pressing on why uniqueness, not mere existence, is the load-bearing claim. One could always
write \emph{some} second-order term and call it a correction; the question is whether two honest
accountants, expanding the same action past first order in two different bookkeeping schemes, would
arrive at the same second-order quantity. The construction proves they would: the remainder is forced,
not chosen. This is the Galileo experiment's demand met at the second order. There it was the count of
recorded events that had to survive every admissible relabeling; here it is the second-order remainder,
and its survival under every way of carving up the expansion is exactly what certifies it as a fact
about the action rather than the accountant. An invariant that two descriptions disagreed on would be a
unit of bookkeeping, struck out as non-physical by the first chapter's rule. The second variation
passes, and its passing is what entitles it to be called the invariant at all.

\section{Its magnitude is the energy}
\label{se:its-magnitude-is-the-energy}

The third road is the one that turns the second variation from a quantity into a measure. The twelfth
chapter built the Sobolev energy, a symmetric positive-definite form that detects zero, and earned its
definiteness at a boundary. The construction now proves that the second variation, read along its
diagonal --- a variation paired with itself --- has a magnitude equal exactly to that energy-squared.
The second variation is not merely \emph{like} the energy; its self-pairing \emph{is} the energy, the
accumulated square of a path's variation. And because the energy was made positive-definite by killing
its nullspace at the boundary, the diagonal second variation inherits that property: it is strictly
positive whenever the activity is nonzero, and zero only when nothing is varying. The second variation
detects zero.

The harmonic oscillator makes the abstraction tangible, and the construction returns to it deliberately
because everything later is built on its bones. Picture the action as a bowl: the stationary path sits
at the bottom, the first variation vanishes there, and the second variation is the curvature of the
bowl's walls. A bowl curves upward in every direction, so every nonzero displacement from the bottom
costs energy --- the second variation is positive in every direction, which is exactly
positive-definiteness, and exactly the statement that the bottom is a genuine minimum. Replace the bowl
with a saddle and the picture breaks: a saddle curves up one way and down another, its second variation
is positive in some directions and negative in others, and a displacement along the downhill direction
costs nothing or less than nothing. The boundary anchor of the previous chapter is what guarantees the
bowl and forbids the saddle --- it kills the flat directions along which a saddle would escape, leaving
a form that is positive everywhere it is not zero. The second variation detects zero precisely because
the construction paid, at the boundary, to make its action a bowl.

There is a reason to dwell on this magnitude, because it is, quite literally, what an instrument reads.
The second variation paired with itself is a single non-negative number --- the construction writes it
as a squared quantity, the activity of a variation totalled and made positive --- and that number is the
reading. When the apparatus of the earlier chapters detects something, what it detects is this: the
magnitude of the second variation at the configuration in front of it, large when there is structure to
register and zero when there is not. The whole edifice of carriers and Facts and decided truth-values
has been, all along, an apparatus for computing one non-negative number and reporting it. The second
variation is not a quantity the construction happens to be able to measure; it is the quantity
measurement \emph{is}, the single reading every instrument in the book returns, and naming it as the
invariant is naming what the book has meant by measurement from its first page.

That is the property that makes it the single invariant the book is named for and not just a term in an
expansion. A measure must report zero for absence and something positive for presence, and the second
variation, by way of its identity with the energy, does exactly that. The cleanest case is the harmonic
oscillator, and it is worth holding onto because every later structure is a variation on it. The
oscillator's action is a quadratic form in a conjugate pair --- a recorded state and the rate at which
that state is expected to change --- and its second variation simply \emph{is} that quadratic form. The
energy of the oscillator, the conserved quantity the stress-energy of the field carries, is the
magnitude of its second variation, positive because the bottom of the oscillator's parabola is a
genuine minimum and not a saddle. The experiment on the harmonic oscillator and the experiment on the
conservation of energy are, read through this chapter, the same statement seen twice: the invariant
that governs the oscillator's motion and the energy that is conserved along it are one number, the
diagonal of the second variation.

\section{The single invariant}
\label{se:the-single-invariant}

Three roads, one destination. The second variation is the mixed coupling carried whole from the
splines; it is the unique quadratic remainder left by the first variation; and its diagonal magnitude
is the positive-definite energy that detects zero. The construction proves all three identifications,
so these are not three quantities that happen to agree but one quantity wearing three descriptions ---
which is the strongest sense in which it is a single invariant. Everything the finite construction has
been doing, from the first cut that made a real to the partitions that approached the continuum up to
epsilon, has been an approach to this one number: the invariant that the finite conditions converge
on, the generic quantity decided in the limit, the thing that survives when every way of describing the
action has been tried and divided away.

Read through the construction's deepest frame, this is the generic the finite conditions were always
approaching. Every cut that made a real, every partition that approximated a path, every refinement
that squeezed a residual toward zero was a finite condition narrowing toward an object it would never
hold in full --- and the second variation is that object, the invariant the conditions decide in the
limit, up to any epsilon one demands but never exactly. The construction tanges it by the one
characteristic that distinguishes a genuine invariant from a bookkeeping artifact: that it is positive
on everything present and zero on nothing, the same under every description, forced rather than chosen.
To possess the invariant up to epsilon is all a finite construction can do and all a measurement ever
needed; the exact value is the limit the conditions approach on a schedule and reach at no finite
stage.

Read one way further, and the single invariant shows a fourth face --- not a fourth thing proved, but the
reading that ties this summit back to the book's first page. The construction's first act was the barest
one available to it: telling two things apart, the first tange, the difference between one Fact and
another, and the order that difference induces. Everything since has been built on it. Now notice what
kind of thing the second variation is --- a difference taken of a difference, a variation of a variation,
the first act turned upon its own result. And recall that the climb did not end by pointing the apparatus
at some configuration in the world; it ended by pointing the apparatus at its own truth, the bare
true-or-false it learned to read of itself at the chapter where the distinction collapsed.
Telling-apart, applied there, has nowhere outward left to point: it reads the first gate against the very
truth that gate decides, and the two readings it would set side by side are forced into one. What remains
is this invariant --- the bare difference, carried the whole way up and standing now over the apparatus's
own true-or-false rather than over any value beneath it. The construction identifies its single invariant
with that self-applied first difference. This is the reading that makes the chapter a summit rather than
merely a high place: the invariant is the first gate of the first chapter turned on itself, and it is
exactly this object the descent will read back down --- as the electron, the first difference returned as
matter.

From here the book changes its reading. Until now each chapter built a piece of apparatus; from now each
chapter reads a value of this one invariant. When a later chapter finds an electron, it will be reading
the second variation of a particular action; when it finds a quantity of boundary radiation, it will be
reading a residue of this same invariant; when it speaks of energy at all, it will mean the diagonal of
this form. The single invariant is the lens, and the remaining chapters are what is seen through it. The
construction set out to measure how true a single Fact is, and it has arrived, by way of a finite
calculus of variations, at the one number that measurement comes down to.

\section{What is demonstrated, and what is assumed}
\label{se:demonstrated-assumed-ch13}

The title theorem deserves the plainest possible accounting of what stands behind it.

What is demonstrated is the threefold identity, proved in full. The second variation is defined as the
mixed coupling and shown to be exactly what remains of the coupled difference once the first-order
residuals vanish. It is proved to be the unique quadratic remainder, so that it is a fact about the
action and not its description. And its diagonal magnitude is proved equal to the Sobolev energy-squared
and strictly positive on nonzero activity, so that it detects zero and serves as a genuine measure.
These are theorems about finite objects, each checked, and together they are the single statement the
book is titled for.

\begin{quote}
\emph{A note on the one assumption.} Nothing in this chapter is granted. The threefold identity is
constructive throughout; no oracle is consulted to prove that the mixed coupling, the unique remainder,
and the energy coincide. The single assumption the book carries remains the one declared at the tenth
chapter --- that stationarity, a universal over infinitely many directions, is decided --- and it stands
in the background here only because clearing the first variation, which let the second variation step
forward, relied on it. The invariant itself is built and proved. What is assumed is only that the path
at which it is read is genuinely the one nature takes, which is the law of motion and no new thing.
\end{quote}

What is assumed, then, is exactly what was assumed two chapters ago, and the second variation is earned
free and clear on top of it. With the single invariant named and shown to be one object reached three
ways, Part II is complete: the construction has its number. The third part turns to what that number
is a number \emph{of} --- to the concrete operators and records that read it, and to the first place a
value of the invariant acquires a physical name of its own.
